\documentclass[12pt]{article}
\usepackage[T1]{fontenc}
\usepackage[utf8]{inputenc}
\usepackage{lmodern}
\usepackage{textcomp}
\usepackage{enumitem}
\usepackage{geometry}
\usepackage{graphicx}
\usepackage{amsmath, amssymb}
\geometry{margin=1in}
\begin{document}
\begin{center}
Economics - K-12 Level Assessment 24 \
Total Marks: 40 \
Answer all questions.
\end{center}

\section*{Multiple Choice (10 marks)}
\begin{enumerate}[label=\arabic*.]
\item Tariffs and quotas
\begin{enumerate}[label=(\alph*)]
    \item result in higher domestic prices.
    \item promote trade between nations.
    \item do not necessarily affect domestic prices.
    \item affect domestic prices: the former raises them while the latter lowers them.
\end{enumerate}
\item To deflate nominal gross domestic product (GDP) you must
\begin{enumerate}[label=(\alph*)]
    \item divide nominal GDP by the GDP deflator.
    \item multiply real GDP by the GDP deflator.
    \item divide real GDP by the GDP deflator.
    \item multiply nominal GDP by the GDP deflator.
\end{enumerate}
\item Expansionary fiscal policy would best be prescribed to
\begin{enumerate}[label=(\alph*)]
    \item eliminate a recessionary gap
    \item reduce inflation
    \item reduce the interest rate
    \item eliminate an inflationary gap
\end{enumerate}
\item If, for each additional unit of a variable input, the increases in output become smaller, which of the following correctly identifies the concept?
\begin{enumerate}[label=(\alph*)]
    \item Diminishing marginal productivity.
    \item Diminishing marginal utility.
    \item Increasing marginal utility.
    \item Increasing marginal productivity.
\end{enumerate}
\item Suppose transfer payments are greater than Social Security contributions corporate taxes and retained earnings combined. In that case
\begin{enumerate}[label=(\alph*)]
    \item NDP will be greater than GDP.
    \item NI will be greater than GDP.
    \item PI will be greater than NI.
    \item DPI will be greater than PI.
\end{enumerate}
\end{enumerate}

\section*{True / False (10 marks)}
\textit{Answer True or False}
\begin{enumerate}[label=\arabic*.]
\setcounter{enumi}{5}
\item True or False: Credit cards are not considered part of the money supply known as M 1.
\item True or False: Expansionary fiscal policy causes less crowding out when government spending improves profit expectations among businesses.
\item True or False: At the natural rate of unemployment, there is no structural unemployment affecting the labor market.
\item True or False: An increase in the price level will shift the aggregate demand curve to the right.
\item True or False: Monopolies can set prices freely without affecting the quantity sold, so their marginal cost curve slopes downward.
\end{enumerate}

\section*{Long Form Response (20 marks)}
\textit{Answer each question with a clear and complete explanation.}
\begin{enumerate}[label=\arabic*.]
\setcounter{enumi}{10}
\item Discuss the concept of monopolistic competition and explain how excess capacity affects firms in such a market structure. Provide examples to support your analysis.
\item Discuss how the purchase of an entertainment system manufactured in China by an American affects U.S. national income accounts, particularly in terms of net exports and GDP.
\end{enumerate}

\vfill
\noindent\textit{End of Paper}
\end{document}